\section{数据理解与总体分析}
\subsection{输入结构与审计}
输入同时提供节点和边的 JSON、CSV 表示。节点编号用于识别对象，不代表最优执行次序。操作节点的缓冲区列表同时包含输入与输出，不能仅依据列表顺序推断读写属性。本文直接使用题面给出的依赖边约束生产与消费先后。

六图规模如下。读取时逐字段比对两种表示，检查节点编号、重复边、缓存类型、大小及申请释放配对。全部输入通过检查，没有对缺失值进行填补，也没有删除任何节点或依赖。
\begin{table}[H]\centering\caption{输入计算图规模}\small
\begin{tabular}{lrr}\toprule 计算图 & 节点数 & 依赖边数\\\midrule
Conv\_Case0 & 2580 & 3869\\
Conv\_Case1 & 36086 & 85653\\
FlashAttention\_Case0 & 1716 & 2712\\
FlashAttention\_Case1 & 6952 & 11184\\
Matmul\_Case0 & 4160 & 7104\\
Matmul\_Case1 & 30976 & 55040\\\bottomrule
\end{tabular}\end{table}

\subsection{操作规模与缓冲区共享}
节点总数包含申请与释放。每个缓冲区恰有一对管理节点，因此操作数满足 $|V_o|=N-2|B|$。表\ref{tab:structure}把真正占用执行单元的操作与缓冲区分别计数，其中“多使用者对象”表示被至少两个操作引用的缓冲区，使用者同时包括生产与消费操作。该指标描述引用范围；判断数据复用还需结合操作属性和依赖关系。
\begin{table}[H]\centering\small
\caption{操作规模与缓冲区引用结构}\label{tab:structure}
\begin{tabular}{lrrrr}\toprule
计算图 & 操作数 & 缓冲区数 & 多使用者对象 & 最多使用者\\\midrule
卷积小图 & 918 & 831 & 823 & 16\\
卷积大图 & 12060 & 12013 & 12013 & 97\\
注意力小图 & 572 & 572 & 540 & 5\\
注意力大图 & 2296 & 2328 & 2200 & 9\\
矩阵乘小图 & 1728 & 1216 & 1216 & 9\\
矩阵乘大图 & 13056 & 8960 & 8960 & 17\\
\bottomrule\end{tabular}
\end{table}

卷积大图有 12060 个操作、12013 个缓冲区，单个缓冲区最多被 97 个操作引用。固定某一操作之后，共享对象仍可能服务于多个后继分支，其释放时刻取决于这些分支中最后完成的使用者。因而用“本操作完成”作为释放依据会丢失后续数据；求解时需要维护每个对象的剩余使用数，并同时检查释放节点的显式前驱。卷积小图的最大使用者数为 16，与大图存在明显差别，后文分别考察两图的驻留组成和重复换出。

矩阵乘小图和大图分别包含 512、4096 个矩阵乘操作，级间移动操作分别为 1024、8192 个。矩阵计算前的两个输入通常由级间搬运提供，因而单纯提前矩阵操作的排序优先级，仍会受到数据到达时刻的约束。两图的缓冲区分别有 1216、8960 个，其中一级缓存对象为 128、512 个，零级左右输入缓存对象合计为 1024、8192 个。对象数量分布与空间压力分布需要分别测量：大量短期零级对象可以分时复用地址，较少而寿命较长的一级对象也可能形成主要驻留峰值。

注意力小图和大图分别包含 64、256 个矩阵乘操作，同时具有行最大值、指数、行求和、乘法与加法等向量操作。其统一缓存对象分别为 356、1512 个，占全部对象的较大部分。这一结构使矩阵单元与向量单元之间的数据交接成为流水分析的重要对象。两个阶段可以由不同硬件单元承担，但后一阶段仍须等待所需数据及缓存地址同时就绪。后文选择注意力小图展示局部时间窗，是因为同一窗口能够观察矩阵计算、级间搬运和向量计算的先后关系。

六图具有不同的操作组成，求解过程仍统一读取节点的执行单元、周期和缓冲区属性。操作名称用于识别管理节点及题面规定的输入复制计费，其余排序与分配由图结构和资源状态确定。这样可以在同一组规则下比较不同计算图，避免将某一图的结果建立在专门编写的顺序上。

\subsection{目标的耦合关系}
缓存驻留峰值低，并不意味着每种缓存都能分配成功。一方面，不同类型不能互相借用空间；另一方面，分散的空闲区间无法容纳较大的连续申请。与此相反，过度追求小峰值可能使原本可重叠的计算串行化。因此首先通过生命周期收缩建立基准，再分别围绕搬运量和时间评价候选。

采用“拓扑顺序—连续分配—依赖重建—流水评价”的计算链。问题一输出原始节点的置换及其峰值，问题二在候选顺序上增加地址和换入换出事件，问题三利用完整方案中的结构依赖重排各单元指令。三问的接口依次增加物理约束，使每一步产生的中间对象都具有明确含义。

同一操作顺序可以对应不同的连续地址方案。即使两种地址方案引起的搬运量相同，地址复用边也可能不同，进而改变跨单元等待。反过来，减小峰值的顺序可能推迟本可提前加载的数据。本文分别计算峰值 $H$、额外搬运量 $D$ 和完成时间 $T$，在问题二按搬运量优先选择，在问题三以该搬运量为上限比较时间。这样的分层评价保留了三问的目标差别，也使无改进的候选能够用于分析目标之间的冲突。

实验以六个原始计算图为共同输入，所有算法采用相同容量、搬运计费与时间递推。方案生成后，从完整调度文本重新构建驻留和流水，核对节点唯一性、原始边、地址互斥及执行时间。验证对象是可执行的完整方案；汇总表据这些方案计算得到。参数比较保持同一输入与评价方式，具体控制条件在第八节给出。
